\documentclass[a4paper,12pt]{article}

\usepackage[utf8]{inputenc}
\usepackage[T1]{fontenc}
\selectlanguage{francais}

%\usepackage{mathptmx}
%\usepackage[scaled]{helvet}
%\usepackage{courier}

\usepackage[margin=2cm,nohead,foot=1.4cm]{geometry}
\usepackage{parskip}
\usepackage{comment}
\usepackage{graphicx}
\usepackage{fancyvrb}

\begin{document}

\begin{flushleft}
\begin{tabular}[t]{l}
\Large{\textbf{Demande d'information}}\\[4mm]
\large{\textbf{relative à la modification des accès au lotissement Zellmatten 2}}\\

\end{tabular}%

	  \thispagestyle{empty}
\noindent%
\begin{tabular}[t]{llr}
	  Stéphane \textsc{Genaud}, & Jean \textsc{Heckly}, &  au nom des résidents du lotissement,\\
	  \small{9 rue des Pommiers} & \small{2 imp. des Cerisiers}\\
	  Walbach            & Walbach.\\
\end{tabular}%
\end{flushleft}

\vspace{6mm}

Monsieur le Maire,
Mesdames et Messieurs les conseilles municipaux,


Le conseil municipal de Walbach s'est réuni le 22 janvier 2019 et a discuté (point 4 de l'ordre du jour)
une modification substantielle (car absente des plans de masse fournis lors de la commercialisation) 
du lotissement Zellmatten 2. Cette modification est l'ouverture aux 
véhicules et piétons du segment actuellement non carossable reliant le lotissement à la 
rue des Acacias.

Le procès verbal du conseil indique que cette discussion s'est conclue par des décisions prises à l'unanimité:
\begin{itemize}
	  \item \emph{Accord pour l’ouverture du passage au Nord du lotissement
Zellmatten 2 à la circulation des piétons et automobilistes, l’aménagement d’un
accès sécurisé devra être étudié ;}
\item \emph{Demande au Maire d’engager toutes les démarches en vue de l’acquisition par
	  la Commune des parcelles concernées.}
\end{itemize}


Il apparaît que l'immense majorité des résidents du lotissement n'ont appris cette décision qu'à travers
un article de presse plus d'un mois plus tard.

Le sentiment général parmi les résidents est double.
\begin{enumerate}
	  \item
Il y a d'abord la surprise : à notre connaissance aucune consultation ni information n'a été délivrée aux 
intéressés. Nous sommes surpris par la méthode. C'est donc le premier point de notre demande d'information:
le conseil municipal va t-il consulter formellement les habitants du quartier, et si oui, sous quelle forme, avant 
d'entreprendre une action définitive ?
	  \item
Le deuxième sentiment est la crainte des conséquences de l'ouverture aux véhicules. Beaucoup s'interrogent
sur la cohérence de sécuriser la rue de la gare pour ensuite ouvrir une nouvelle voie qui pourrait s'y substituer.
La décision prise en conseil a été d'ouvrir, en relayant l'étude de la sécurité à une échéance ultérieure non précisée.
Nous sollicitons donc une information sur ce deuxième point : 
\begin{itemize}
	  \item quels sont les arguments qui amènent  le conseil
		municipal à se prononcer à l'unanimité pour cette ouverture ?
	\item pourquoi la décision d'ouverture n'est t-elle pas prise sur la base d'un projet précis incluant les 
		mesures de sécurité prévues (et donc chiffrées financièrement).
\end{itemize}
\end{enumerate}

Nous serions donc reconnaissants à la municipalité de nous apporter toutes les informations sur le projet, et en 
particulier les accords trouvés avec le lotisseur, afin d'éclairer le débat.

Nous vous prions, Mesdames, Messieurs les élus, d'agréer nos sentiments respectueux.


\begin{flushright}
{\small Fait à Walbach, le 3 mars 2019}.\\
%\includegraphics[width=.16\textwidth]{signgenaud.jpg}
\end{flushright}
\end{document}
